% Integrated after independent root, criteria, and geometric audits, 2026-11-16.
% Original author draft retained in agent_F/F_fixed_r_boundary_draft.tex.
\subsection{The diagonal boundary and transport of cuspidal groups}

Retain the arithmetic notation of
\ref{clm:AGfixedrconfiguration}: \(C/k\) is smooth proper
of genus at least two over a finite field of characteristic
\(p\), \(r\ne p\), and
\(\mathcal E_{\mathrm{c}r}\) is the cuspidally pro-\(r\)
quotient of the geometrically prime-to-\(p\) group of
\(U=(C\times C)-\Delta_C\), relative to
\(\Pi_C\times_{G_k}\Pi_C\).

\begin{claim}{AG.42}{Proved: diagonal boundary and cuspidal transport}
\label{clm:AGcuspidaltransport}
There is a diagonal boundary decomposition subgroup
\(\mathcal D\subset\mathcal E_{\mathrm{c}r}\) with
an exact sequence
\[
 1\longrightarrow\mathbb Z_r(1)\longrightarrow\mathcal D
 \longrightarrow\Pi_C\longrightarrow1,
                                                        \tag{AG.42.1}
\]
whose map to the proper product is the diagonal map on
the last term. For every rational point \(P\), choose
compatible geometric paths through this boundary and put
\(D_P=s_P(G_k)\subset\Pi_C\) and
\(\mathcal E_P=\mathrm{pr}_2^{-1}(D_P)\).
The fixed-\(r\) configuration theorem identifies
\(\mathcal E_P\) with the arithmetic cuspidally
pro-\(r\) group of \(C-P\), and under this
identification its puncture decomposition group is
\[
                           \mathcal D\cap\mathcal E_P.
                                                        \tag{AG.42.2}
\]

If two rational points \(P,Q\) have geometrically conjugate
proper sections in \(\Pi_C\), there is an isomorphism
between the actual cuspidally pro-\(r\) groups of
\(C-P\) and \(C-Q\), over the identity of
\(\Pi_C\), which carries the first puncture decomposition
group to a conjugate of the second. Consequently
\(P=Q\): rational point sections of the full proper
geometrically prime-to-\(p\) fundamental group are
injective up to geometric conjugacy.
\end{claim}

\begin{proof}
\emph{Construction of the boundary group.}
We give the local tame-cover construction explicitly to
isolate it from any full configuration-fiber assertion.
Let \(X=C\times C\), \(D=\Delta_C\), and
\(N=\mathcal O_X(D)|_D\), the normal line bundle.
For each positive integer \(n\) prime to \(p\), form
the root stack \(X_n=\sqrt[n]{(X,D)}\). Locally where
\(D\) has parameter \(t\), this is the quotient of
\(\operatorname{Spec}(A[z]/(z^n-t))\) by \(\mu_n\).
The reduced inverse image \(D_n\) of \(D\) is the
\(\mu_n\)-gerbe of \(n\)-th roots of \(N\).

A finite etale cover of \(U\) with prime-to-\(p\)
geometric monodromy has tame inertia along \(D\).
Take \(n\) divisible by its finitely many inertia
orders. The absolute Abhyankar lemma, SGA~1 XIII,
Proposition~5.2, says that after the local substitution
\(z^n=t\), its normalization extends uniquely to a
finite etale cover of the whole Kummer chart
\cite{AR_SGA1}. The cited statement already includes
extension over the entire chart; its proof applies
purity after the codimension-one tame calculation.
The unique extensions and their deck
descent data glue to a finite etale cover of \(X_n\).
Restriction to \(D_n\) gives a finite etale cover of
that gerbe. All constructions commute with increasing
\(n\), with field extension, and with pullback
transverse to \(D\). They therefore define the boundary
homomorphism to the tame fundamental group of \(U\),
and hence to its geometrically prime-to-\(p\) quotient.

Here are the exact group assertions needed from this
construction. The fundamental group of the gerbe \(D_n\)
maps onto \(\pi_1(D)\), and its kernel is the image
of \(\mu_n\). This is the elementary gerbe exact
sequence: etale locally on \(D\) the gerbe is
\(B\mu_n\); a cover on which the stabilizer acts
trivially descends to \(D\), and covers of \(D\)
pull back fully faithfully. The inertia is normal,
with the cyclotomic descent action. No injection of
\(\mu_n\) is required here. Passing to compatible
root orders gives the tame boundary group \(\mathcal T\)
and a right exact sequence
\[
 \widehat{\mathbb Z}^{(p')}(1)\longrightarrow\mathcal T
              \longrightarrow\pi_1(D)\longrightarrow1.
                                                        \tag{AG.42.3}
\]
The exactness follows either from the same descent
criterion or by inverse limits of the compact gerbe
sequences; the transition maps are surjective, as is
seen in the local maps \(\mu_m\to\mu_n\) for
\(n\mid m\).

Now take the maximal geometrically prime-to-\(p\)
quotients. On geometric groups, right exactness of
completion shows that quotienting the completed
boundary group by the image of inertia gives
\(\Delta_C\). On arithmetic groups the quotients
over \(G_k\) are unchanged. Thus
\[
 \widehat{\mathbb Z}^{(p')}(1)\longrightarrow
 \mathcal T^{\mathrm{geom},p'}\longrightarrow
 \Pi_C\longrightarrow1                             \tag{AG.42.4}
\]
is right exact. This step explicitly rules out an extra
kernel arising from the discarded geometric base kernel.

Let \(\mathcal D\) be the image of
\(\mathcal T^{\mathrm{geom},p'}\) in
\(\mathcal E_{\mathrm{c}r}\). Its projection to the
proper product factors through the injective diagonal
\(\Pi_C\to\Pi_C\times_{G_k}\Pi_C\).
Consequently the kernel of \(\mathcal D\to\Pi_C\)
is exactly the image of the inertia in (AG.42.4).
Since it is contained in the cuspidal pro-\(r\) kernel,
this inertia image is a quotient of \(\mathbb Z_r(1)\).

This quotient is the full \(\mathbb Z_r(1)\).
To verify this without using a global boundary theorem,
take a geometric fiber \(C-\{P\}\). Its cuspidal
pro-\(r\) group maps onto its maximal pro-\(r\)
fundamental group, which is free on the usual
\(2g\) surface generators. Its puncture inertia
generator is, up to inversion, the product
\(\prod_{i=1}^g[a_i,b_i]\). Map this free group to
the pro-\(r\) upper-unitriangular Heisenberg group
over \(\mathbb Z_r\) by
\[
 a_1\longmapsto 1+e_{12},\qquad
 b_1\longmapsto 1+e_{23},\qquad
 a_i,b_i\longmapsto1\quad(i>1).
                                                        \tag{AG.42.5}
\]
The puncture generator maps to \(1+e_{13}\), of
infinite \(r\)-adic order. Thus its \(r\)-inertia
is \(\mathbb Z_r(1)\). The fixed-\(r\) fiber
injection carries this inertia injectively into
\(\mathcal E_{\mathrm{c}r}\). In the local Kummer
charts the fiber meets the diagonal transversely
with multiplicity one, so this is exactly the
diagonal inertia in (AG.42.4), with its ordinary
Kummer generator. This proves (AG.42.1).

Finally, the same transverse local chart sends the
actual puncture decomposition group of \(C-P\)
into \(\mathcal D\). It lies over \(D_P\), and
it maps isomorphically on inertia and on \(D_P\).
The inverse image of \(D_P\) in (AG.42.1) is an
extension of those same two groups. The resulting
map is therefore an isomorphism onto that inverse
image, proving (AG.42.2). Compatible paths can be chosen
simultaneously using the common boundary gerbes;
other choices give the permitted geometric inner
conjugacies.

\emph{Transport from conjugate proper point sections.}
Suppose \(P,Q\in C(k)\) have geometrically conjugate
proper sections. Choose
\(\gamma\in\Delta_C\) with
\(\gamma D_P\gamma^{-1}=D_Q\).
By (AG.42.1), choose a geometric
\(d\in\mathcal D\) mapping to \(\gamma\).
Conjugation by \(d\) identifies
\(\mathcal E_P\) with \(\mathcal E_Q\), and by
(AG.42.2) carries their puncture decomposition groups
onto each other. On the first proper coordinate
it induces \(\operatorname{Inn}(\gamma)\).

The first projection
\(\mathcal E_Q\to\Pi_C\) is surjective, and is
surjective on geometric kernels: it is the proper
quotient of the punctured curve group, by the
fixed-\(r\) theorem. Choose a geometric
\(e\in\mathcal E_Q\) mapping to \(\gamma^{-1}\)
under that first projection. Then
\[
            \beta=\operatorname{Inn}(e)\circ
                         \operatorname{Inn}(d):
                 \mathcal E_P\xrightarrow{\sim}\mathcal E_Q
                                                        \tag{AG.42.6}
\]
is over the identity of \(\Pi_C\), is over the
identity of \(G_k\), and carries the first puncture
decomposition group to a conjugate of the second
inside \(\mathcal E_Q\). These are exactly the
hypotheses of Claim~\ref{clm:AGcuspidalcriterion}.
The construction uses the one-prime relative fiber
comparison of Claim~\ref{clm:AGfixedrconfiguration}
and the tame boundary comparison just proved; it does
not invoke full prime-to-\(p\) configuration fiber
exactness. Claim~\ref{clm:AGcuspidalcriterion} now
implies \(P=Q\), as asserted. The argument applies
unchanged after every finite extension of the constant
field.
\end{proof}
